\begin{abstract}
	This paper introduces the Nosy Neighbor, a Kubernetes anti-pattern where a cluster-scoped Operator watches and caches resources across all namespaces despite only needing a small subset.
	Driven by clean development clusters and default controller-runtime caching behaviors, this issue often goes unnoticed until it causes production OOMKilled crashes.
	To prevent and remediate the problem, this paper outlines practical steps for auditing RBAC resource classifications, filtering caches via resource labels, and running synthetic load tests to catch memory spikes early.
\end{abstract}
